
\multiproblem{linear}{

    We now consider the equation
    \begin{equation}
        a\,x = b
        \label{eq:axb}
    \end{equation}
    where $a$ and $b$ are known and $x$ is an unknown quantity. A solution of
    the equation is a value for $x$ that satisfies the equation. Our objective
    is to find the set of all possible solutions to this equation given
    different $a$ and $b$.
    \begin{enumerate}
        \item Let's first consider some examples. What is the solution set
            for $x$ in the following cases?
            \begin{enumerate}
                \item $2\,x = 8$
                \item $0\,x = 8$
                \item $0\,x = 0$
            \end{enumerate}
        \item More generally given $a,b,x\in\mathbb{R}$ how does the solution
            set for~\eqref{eq:axb} depend on $a$ and $b$?
        \item Given any particular values for $a$ and $b$ can you say how many
          different solutions for $x$ there will be?
        \item A linear equation can only ever have $0$, $1$, $\infty$
          solutions. Can you think of types of equations that can have other
          numbers of solutions (e.g. $2$ solutions)?
    \end{enumerate}
}


\multiproblem{lineareqns}{
    Solve the following systems of equations using Gaussian elimination. Note
    that the first example is easier to solve by substitution - the advantage
    of Gaussian elimination becomes clearer as the system gets larger.
    \begin{enumerate}
        \item
            \begin{align*}
                x + 7y &= -13 \\
                2x + y &= 0
            \end{align*}
        \item
            \begin{align*}
                -x + 3y + 4z &= 5  \\
                7x +  y + 6z &= -3 \\
                2x -  y -  z &= 3  \\
            \end{align*}
        \item
            \begin{align*}
                2x + 6y - 2z - 2w &= -9 \\
                4x + 12y +2z - 2w &= 4  \\
                     2y - 6z + 2w &= -13\\
                2x + 2y + 2z + 2w &= 9
            \end{align*}
    \end{enumerate}
}

\multiproblem{nonunique}{
    In general systems of linear equations can be written in matrix form
    \[
        \mathbf{Mx} = \mathbf{y}
    \]
    where $\mathbf{M}$ is a (known) matrix, $\mathbf{y}$ is a known vector and
    $\mathbf{x}$ is an unknown vector e.g.
    \[
        \mathbf{x} =
            \begin{pmatrix}
                x \\ y \\ z
            \end{pmatrix}
    \]
    Given this equation in matrix form it is clear that if the matrix is
    invertible then we have a unique solution since we can multiply both sides
    to find that
    \[
        \mathbf{M}^{-1} (\mathbf{Mx}) = \mathbf{M}^{-1}\mathbf{y}
    \]
    Working on the left hand side we see that (using associativity of matrix
    multiplication and definition of the inverse matrix)
    \[
        \mathbf{M}^{-1} (\mathbf{Mx}) =
        (\mathbf{M}^{-1} \mathbf{M})\mathbf{x} =
        \mathbf{I}\mathbf{x} =
        \mathbf{x}.
    \]
    Putting this together we have
    \[
        \mathbf{x} = \mathbf{M}^{-1}\mathbf{y}.
    \]
    So if the matrix $\mathbf{M}$ is invertible then there is a unique
    solution to the system of equations. The matrix is invertible if and only
    if $\mathrm{det}\,\mathbf{M}\neq 0$. Otherwise the system of equations may
    have zero solutions or infinitely many solutions. In each of the following
    problems you should calculate the determinant \emph{and} use Gaussian
    elimination to see what happens.

    \begin{enumerate}
        \item Find the set of solutions to
            \begin{align*}
                x + y + z  &= 1 \\
                x - y - 2z &= 2 \\
                4x - 2z    &= 6
            \end{align*}
            (The set of solutions is of the form $\{\mathbf{v}_1 + \alpha
            \mathbf{v}_2 : \alpha \in \mathbb{R}\}$ for some vectors
            $\mathbf{v}_1$ and $\mathbf{v}_2$ that you should find.)
        \item Find the set of solutions to
            \begin{align*}
                x + y + z  &= 1 \\
                x - y - 2z &= 2 \\
                4x - 2z    &= 5
            \end{align*}
        \item For what values of $c$ does the following system \emph{not} have
            a unique solution?
            \begin{align*}
                x  -  y +  z &= 1 \\
                2x + cy + 3z &= 0 \\
                      y + 2z &= -1
            \end{align*}
            What is the set of solutions in that case?
    \end{enumerate}
}

\multiproblem{polynomials}{
    Consider trying to find the coefficients of a polynomial $P(x)$ whose
    curve $y=P(x)$ passes through a set of points in the $x,y$ plane. This
    naturally gives us a system of linear equations that may or may not have a
    unique solution.
    \begin{enumerate}
        \item We want to find a polynomial that passes through the points
            $(-1,8)$, $(0, 3)$ and $(1, 0)$. Find the set of all quadratic
            (i.e. $P(x)=ax^2+bx+c$) polynomials that pass through these three
            points.
        \item Find the set of all linear polynomials ($P(x)=ax+b$) passing
            through the same three points (draw a diagram to interpret this
            result).
        \item Find the set of all quadratic polynomials passing through only the
            two points $(0, 3)$ and $(1,0)$ - sketch a few of them.
        \item Consider trying to find a unique quadratic polynomial passing
            through three points $(x_1,y_1)$, $(x_2,y_2)$ and $(x_3,y_3)$.
            Under what conditions on the coordinates $x_i$ and $y_i$ will
            there be a unique polynomial?
    \end{enumerate}
}

\multiproblem{supports}{
    Here we consider problems involving force and moment balance which also
    naturally lead to systems of linear equations\\
    \begin{minipage}{.65\textwidth}
    \begin{enumerate}
        \item A rectangular body of length $4\,\mathrm{m}$ and height
            $2\,\mathrm{m}$ is in equilibrium subject to three different
            forces. The weight $mg$ acts at the centre of the body. Another
            force $F_1$ acts at the upper left corner at an angle of
            $30^{\circ}$ to the horizontal. A third force $F_2$ acts
            $1\,\mathrm{m}$ from the upper-right corner at an angle $\alpha$.
            Find the magnitude of the unknown force $F_1$ and also the force
            vector for $F_2$. Hence find the angle $\alpha$. Show that the
            lines of action of all three forces intersect at a single point
            (and find that point).\\
        \item Consider a rod of (known) length $l$ and mass $m$ resting on two
            supports as shown
            Using force and moment balance find the two (unknown) support
            forces $N_1$ and $N_2$ and verfiy that they satisfy the required
            constraint for a resting support ($N_i \ge 0$).\\
        \item Consider the same but with three supports
            Using force and moment balance find the set of possible values for
            the three support forces $N_1$, $N_2$, and $N_3$. Can we get a
            unique solution by taking moments about more than one place?
            Considering the resting constraint ($N_i \ge 0$) find minimum and
            maximum values for each of the three support forces.
    \end{enumerate}
    \end{minipage}
    \begin{minipage}{0.3\textwidth}
      \begin{center}
        \includegraphics{moments.pdf}\\
        \vspace{.2in}
        \includegraphics{supports2.pdf}\\
        \vspace{.2in}
        \includegraphics{supports3.pdf}\\
      \end{center}
    \end{minipage}
}

\multiproblem{algebraic}{
    Given a system of linear equationsin matrix form as
        $\mathbf{Mx} = \mathbf{y}$
    if $\mathbf{y}=\mathbf{0}$ then the system takes the form
    \begin{equation}
      \mathbf{Mx} = \mathbf{0}
      \label{eq:homo}
    \end{equation}
    and we say that the system is \emph{homogeneous}.
    \begin{enumerate}
      \item Show that $\mathbf{x}=\mathbf{0}$ is always a solution of a
        homogeneous linear system (regardless of what $\mathbf{M}$ is).
      \item Show that~\eqref{eq:homo} will have a non-zero solution for
        $\mathbf{x}$ if and only if $\mathrm{det}\,\mathbf{M}=0$.
        \item Show that if some vector $\mathbf{x}_1$ is a solution
          of~\eqref{eq:homo} then any scalar multiple of $\mathbf{x_1}$ (e.g.
          $2\mathbf{x}_1$) is also a solution.
        \item Show that given two solutions $\mathbf{x_1}$ and $\mathbf{x_2}$
          of~\eqref{eq:homo} $a\,\mathbf{x}_1 + b\,\mathbf{x}_2$ is a solution
          for any scalars $a$ and $b$.
        \item Which of the previous results also holds for
            \emph{non-homogeneous} systems of equations (with
            $\mathbf{y}\neq \mathbf{0}$.)?
    \end{enumerate}
}
